\documentclass[xcolor=dvipsnames]{beamer}
\usetheme[height=7mm]{Rochester}
\setbeamertemplate{items}[ball]
\setbeamertemplate{blocks}[rounded][shadow=true]
\usepackage{xcolor}
\usepackage{colortbl}
\usepackage{multirow}
\usepackage{amsmath}
\usepackage{amsfonts}
\usepackage{epsfig}
\usepackage{times}
\usepackage[T1]{fontenc}
\usepackage{inconsolata}
\usepackage{fancyvrb}
\usepackage{bbm}

% == theme & colors
\usetheme{Warsaw}
\definecolor{someblue}{rgb}{0,.1,0.8}

\title{Stat 256: Causality}
\date{Spring 2026}
\subtitle[ ]{Introduction}
\author{\textbf{Instructor}: Chad Hazlett}
\institute[UCLA]{UCLA\\[1ex]}

\definecolor{dgreen}{rgb}{0,.6,0.8}
\newtheorem{assumption}{Assumption}
\newtheorem{iass}{Identification Assumption}
\newtheorem{ires}{Identification Result}
\newtheorem{estm}{Estimand}
\newtheorem{esti}{Estimator}
\newtheorem{question}{Question}
\newtheorem{answer}{Answer}
\newtheorem{proposition}{Proposition}
\newcommand{\bs}{\boldsymbol}
\newcommand{\mb}{\mathbf}
\newcommand{\real}{\mathbb{R}}
\renewcommand{\u}{\ensuremath{\mb{u}}}
\renewcommand{\v}{\ensuremath{\mb{v}}}
\newcommand{\U}{\ensuremath{\mb{U}}}
\newcommand{\Z}{\ensuremath{\mb{Z}}}
\newcommand{\PP}{\ensuremath{\mb{P}}}
\newcommand{\M}{\ensuremath{\mb{M}}}
\newcommand{\X}{\ensuremath{\mb{X}}}
\newcommand{\x}{\ensuremath{\mb{x}}}
\newcommand{\Y}{\ensuremath{\mb{Y}}}
\newcommand{\D}{\ensuremath{\mb{D}}}
\newcommand{\y}{\ensuremath{\mb{y}}}
\renewcommand{\b}{\ensuremath{\bs{\beta}}}
\newcommand{\e}{\ensuremath{\bs{\epsilon}}}
\newcommand{\bhat}{\ensuremath{\hat{\bs{\beta}}}}
\newcommand{\XX}{\ensuremath{\X'\X}}
\newcommand{\XXinv}{\ensuremath{\left(\XX\right)^{-1}}}
\newcommand{\hatsig}{\ensuremath{\hat{\sigma}^2}}
\newcommand{\sig}{\ensuremath{\sigma^2}}
\newcommand{\hatmat}{\ensuremath{\X \XXinv \X'}}
\newcommand{\ehat}{\ensuremath{\hat{\e}}}
\newcommand{\tr}{\text{tr}}
\newcommand{\Sig}{\ensuremath{\bs{\Sigma}}}
\newcommand{\SigInv}{\ensuremath{\bs{\Sigma^{-1}}}}
\newcommand{\W}{\ensuremath{\mb{W}}}
\newcommand\E{\mathbb{E}}

\newcommand{\indep}{{\bot\negthickspace\negthickspace\bot}}
\useoutertheme{infolines}
\begin{document}

\begin{frame}
  \titlepage
\end{frame}

%% ============================================================
%% WHY CAUSALITY
%% ============================================================

\begin{frame}{Why causality?}
\onslide<1-> \textit{Most} interesting questions in science and life are causal. \\
\footnotesize
\begin{itemize}
\item Does this drug / policy / event cause higher or lower $Y$?
\item What explains why $Y$ happened?
\end{itemize}
\bigskip
\normalsize
\onslide<2-> And yet:\\
\footnotesize
\begin{itemize}
\item<3-> There is no causal inference without assumptions that are fundamentally unverifiable from data alone. \medskip
\item<4-> Traditional statistics does not include a coherent language for causal effects, much less tools. \medskip
\item<5-> Even ``descriptive'' analyses involve missingness, case selection, and generalization---problems rooted in causal processes. \medskip
\item<6-> Knowing what you can and cannot say, and protecting against abuse of your results, is vital to being an honest statistician.
\end{itemize}
\end{frame}

%% ============================================================
%% GOALS
%% ============================================================

\begin{frame}{Goals}
\begin{itemize}
\item<1-> Equip you with tools for careful causal thinking \medskip
\item<2-> Equip you to read, critique, and (prepare to) write methodological work on causal inference \medskip
\item<3-> Provide you with tools for applied research and the wisdom to
\begin{itemize}
\item avoid abusing these tools
\item carefully assess the credibility of claims made by others
\end{itemize}
\medskip
\item<4-> What we focus \textit{less} on:
\begin{itemize}
\item some of the more esoteric theoretical directions
\item estimation horse-races while deeper problems go ignored
\end{itemize}
\end{itemize}
\end{frame}

%% ============================================================
%% THEMES
%% ============================================================

\begin{frame}{Theme 1: Bilingualism}
\onslide<2-> Two frameworks dominate work on causal inference:\\
\medskip
\onslide<3-> 1. \textbf{Structural causal models (SCMs) / directed acyclic graphs (DAGs)}
\small
\begin{itemize}
\item Encode which variables cause which; derive what you can learn from observational data.
\item Largely attributed to UCLA's own Judea Pearl.
\end{itemize}
\bigskip
\normalsize
\onslide<4-> 2. \textbf{Potential outcomes (POs)}
\small
\begin{itemize}
\item Each unit has potential outcomes under each treatment; write these down and make assumptions about how they relate to treatment assignment.
\item Associated with Jerzy Neyman and Don Rubin.
\end{itemize}

\normalsize
\begin{center}\onslide<5-> These have sometimes been treated as rival camps, but that need not be the case. \textit{We will develop fluency in both.}
\end{center}
\end{frame}

\begin{frame}{Theme 2: Feasibility}
\small
\begin{itemize}
\item<1-> Great theoretical progress has been made in causal inference.
\medskip
\item<2-> But in theory you can assert an assumption whether it is credible or not.
\medskip
\item<3-> We care about how investigators can \textit{make or defend} causal claims in practice.
\medskip
\item<4-> ``Identification strategies'' link causal quantities to observable ones. For each, we ask:
\footnotesize
\begin{itemize}
\item What do the assumptions mean? What about nature or the research design supports them? How could you falsify them? \\
\item How do we avoid over-confidence built on unreliable assumptions? \\
\item How do we make progress when sharp assumptions aren't defensible?
\end{itemize}
\end{itemize}
\end{frame}

%% ============================================================
%% ROAD MAP
%% ============================================================

\begin{frame}{Road map: Part I --- Core language and theory}
\textbf{Unit 1: DAGs and structural causal models}\\
\small
\begin{itemize}
	\item Graphs, d-separation, back-door and front-door criteria, do-calculus
	\item Where do potential outcomes come from?
\end{itemize}
\medskip

\normalsize
\textbf{Unit 2: Potential outcomes}\\
\small
\begin{itemize}
\item Definitions, estimands, randomized experiments
\item Conditional ignorability / selection on observables
\item Bridging POs and DAGs
\end{itemize}
\end{frame}

\begin{frame}{Road map: Part II --- Identification strategies}
\small
\textbf{Unit 3: Selection on observables in practice}\\
\footnotesize
\begin{itemize}
\item Sub-classification, matching, regression, propensity scores, balancing weights
\item Brief survey of double-robust / ML-augmented approaches
\item Sensitivity analysis for SOO
\end{itemize}
\medskip
\small
\textbf{Unit 4: Instrumental variables}\\
\footnotesize
\begin{itemize}
\item IV in PO and DAG frameworks
\item Practical applications and critiques, sensitivity analysis
\end{itemize}
\medskip
\small
\textbf{Unit 5: Regression discontinuity}\\
\footnotesize
\begin{itemize}
\item Sharp and fuzzy RD
\end{itemize}
\medskip
\small
\textbf{Unit 6: Difference-in-differences}\\
\footnotesize
\begin{itemize}
\item Parallel trends; DID $\neq$ TWFE; staggered adoption; sensitivity
\end{itemize}
\end{frame}

\begin{frame}{Road map: Part II (continued)}
\small
\textbf{Unit 7: Time-varying confounding / panel methods}\\
\footnotesize
\begin{itemize}
\item Synthetic control, LDV, interactive fixed effects, relaxing parallel trends
\end{itemize}
\medskip
\small
\textbf{Unit 8: Sensitivity analysis and partial identification}\\
\footnotesize
\begin{itemize}
\item OVB for bounding; sensitivity for IV and other strategies
\item Placebo treatments and outcomes; stability-controlled quasi-experiments
\end{itemize}
\medskip
\small
\textbf{Unit 9: Mediation and TRACE}\\
\footnotesize
\begin{itemize}
\item Why mediation is fundamentally harder; direct and indirect effects
\item Marginal structural models; separating units rather than paths
\item The TRACE: separating units rather than effects to address post-treatment problems
\end{itemize}
\end{frame}

\begin{frame}{Road map: Part III --- You teach!}
\small
\begin{itemize}
\item Groups of 2--3 choose an advanced topic
\item 25-minute lecture + 5 min Q\&A
\item Submit a tutorial with worked examples and problems
\end{itemize}

\onslide<2-> Example topics:\\
\footnotesize
\begin{itemize}
\item Causal inference in networks / interference
\item Causal attribution, necessity, sufficiency (causes of effects)
\item Generalization, transportability, external validity
\item Advanced heterogeneous effect estimation (CATE, causal forests)
\item Negative outcome controls, placebos, falsification
\item Other ID strategies not covered in class
\item \textit{Custom topics} (by permission)
\end{itemize}
\end{frame}

\begin{frame}{More on presentations}
\small
\onslide<1-> Why do this?\\
\footnotesize
\begin{itemize}
\item Mimics the experience of a smaller advanced seminar
\item Learning something well enough to present it produces far better understanding
\item You are well situated to explain new material to your peers
\end{itemize}
\small
\medskip
\onslide<2-> Common pitfalls: \\
\footnotesize
\begin{itemize}
\item Not enough attention to what the identification assumptions really are, how to understand them, what it would take to believe them.
\item Splitting up the topic so that nobody understands the big picture.
\end{itemize}
\end{frame}

%% ============================================================
%% LOGISTICS
%% ============================================================

\begin{frame}{Requirements and grading}
\small
\begin{itemize}
\item<1-> \textbf{Problem sets (25\%)}: Expect 4--5 over the quarter. Posted and submitted online. Ask for extensions \textit{before} the deadline.
\medskip
\item<2-> \textbf{Advanced topic presentation (25\%)}: Groups of 2--3; lecture + tutorial.
\medskip
\item<3-> \textbf{Critical review memo (25\%)}: Individual. Either critique an empirical paper's identification strategy, or take a deep dive into a methods paper (simulations, failure modes, extensions). Due Week 10.
\medskip
\item<4-> \textbf{Active-learning activities (25\%)}: Near-daily in-class exercises to rehearse and reinforce the material. Graded lightly (0/1/2/3). Lowest two dropped. No make-ups.
\end{itemize}
\end{frame}

\begin{frame}{Timeline (preliminary)}
\begin{itemize}
\item \textbf{Week 3}: Sign up for presentation groups and topics.
\medskip
\item \textbf{Week 6}: Check-in meeting with instructor on presentation and memo.
\medskip
\item \textbf{Weeks 9--10}: Student-led lectures.
\medskip
\item \textbf{Week 10}: Critical review memo due.
\medskip
\item \textbf{Finals week}: All presentation materials due.
\end{itemize}
\end{frame}

\begin{frame}{Prerequisites}
We want to accommodate various fields and training backgrounds, but the skills we count on are:
\medskip
\small
\begin{itemize}
\item Essential: probability theory, statistics including inference and regression.
\item Useful: linear algebra, multivariate calculus.
\item Essential: enough \texttt{R} for AI to help you, and a way to write in \TeX{} (e.g.\ Rmarkdown).
\end{itemize}
\end{frame}

\begin{frame}{Books}
\small
\onslide<1-> We will share references and reading expectations as needed.  \\
\bigskip
\onslide<2-> For Unit 1 on DAGs (starting Week 1) you should have:
\begin{itemize}
\item Pearl, J., Glymour, M., \& Jewell, N. P. (2016). \textit{Causal Inference in Statistics: A Primer}. Wiley.
\item Hern\'an, M.A.\ \& Robins, J.M.\ (2024). \textit{Causal Inference: What If}. Chapman \& Hall/CRC. (Free online)
\end{itemize}
\bigskip
\onslide<3-> Other useful references:
\begin{itemize}
\item Pearl, J. (2009). \textit{Causality}. Cambridge University Press. 2nd ed.
\item Cunningham, S. (2021). \textit{Causal Inference: The Mixtape}. (Free online)
\item Huntington-Klein, N. (2022). \textit{The Effect}. (Free online)
\end{itemize}
\end{frame}

\begin{frame}{AI tools}
\small
You are welcome to use AI tools (ChatGPT, Claude, Copilot, etc.) as learning aids.\\
\medskip
\begin{itemize}
\item Problem set write-ups must reflect \textit{your own} understanding.
\item Your critical review memo must be your own analytical work---but discussing ideas with AI to develop your thinking is fine and encouraged.
\item If you use AI substantively, briefly note how.
\end{itemize}
\medskip
The goal is to learn the material. AI can help with that if you use it as leverage rather than a way out.
\end{frame}

\begin{frame}{Other logistics}
\small
\begin{itemize}
\item Lectures: Monday, Wednesday, 11:00--12:15, Haines A44. In-person only. \medskip
\item Slides will be posted on the course website.\medskip
\item My office hours: Wednesday 3:30--4:45, Bunche 3264  \medskip
\end{itemize}
\end{frame}

\begin{frame}{}
\begin{center}
\Large Questions?
\end{center}
\end{frame}

\begin{frame}{Introductions}
\begin{itemize}
\item Your name
\item Your department / program / year
\item What brings you to this course?
\item One causal question you find interesting (any domain)
\end{itemize}
\end{frame}

\end{document}
